\chapter{واژه‌های زبان سلام}
\label{app:keywords}

واژه‌هایی که در این پیوست آمده‌اند، برای خود زبان سلام رزرو شده‌اند و
نمی‌توان از آن‌ها به عنوان نام متغیر یا کارکرد استفاده کرد؛ نه به تنهایی و
نه به عنوان بخشی از یک نام چندواژه‌ای. بخش نخست واژه‌هایی است که در این
کتاب به کار رفته‌اند و بخش دوم واژه‌هایی که در کتاب‌های بعدی با آن‌ها آشنا
می‌شوید.

\section{واژه‌های به‌کاررفته در این کتاب}

\begin{longtable}{lp{8.5cm}l}
\toprule
واژه & معنی و کاربرد & فصل \\
\midrule
\endhead
\fcode{کارکرد} & تعریف کارکرد؛ \fcode{کارکرد آغازین} نقطه‌ی شروع برنامه & \ref{ch:first}، \ref{ch:functions} \\
\fcode{آغازین} & نام کارکرد آغازین برنامه & \ref{ch:first} \\
\fcode{پایان} & پایان هر بلوک (کارکرد، شرط، حلقه) & \ref{ch:first} \\
\fcode{چاپ} & نوشتن روی خروجی و رفتن به خط بعد & \ref{ch:first} \\
\fcode{بنویس} & نوشتن روی خروجی بدون رفتن به خط بعد & \ref{ch:first} \\
\fcode{گذرا} & تعریف متغیری که مقدارش عوض می‌شود & \ref{ch:variables} \\
\fcode{پایدار} & تعریف یک پایدار & \ref{ch:variables} \\
\fcode{درست} & مقدار منطقی درست & \ref{ch:types} \\
\fcode{نادرست} & مقدار منطقی نادرست & \ref{ch:types} \\
\fcode{برگردان} & تبدیل ریخت & \ref{ch:types} \\
\fcode{و} & «و» منطقی & \ref{ch:operators} \\
\fcode{یا} & «یا» منطقی & \ref{ch:operators} \\
\fcode{اگر} & شرط & \ref{ch:conditions} \\
\fcode{وگرنه} & راه دیگر شرط؛ با شرط یا بی‌شرط & \ref{ch:conditions} \\
\fcode{برگزین} & مقایسه‌ی یک مقدار با چند حالت & \ref{ch:conditions} \\
\fcode{چرخه} & حلقه با تعداد مشخص یا بازه & \ref{ch:loops} \\
\fcode{با} & نام‌گذاری شمارنده‌ی حلقه & \ref{ch:loops} \\
\fcode{تا} & پایان بازه در \fcode{چرخه} & \ref{ch:loops} \\
\fcode{گام} & اندازه‌ی پرش در \fcode{چرخه} & \ref{ch:loops} \\
\fcode{تاهنگام} & حلقه تا وقتی شرط برقرار است & \ref{ch:loops} \\
\fcode{بشکن} & پایان دادن به حلقه & \ref{ch:loops} \\
\fcode{بگذر} & پریدن به دور بعدی حلقه & \ref{ch:loops} \\
\fcode{بازگشت} & پس دادن نتیجه از کارکرد & \ref{ch:functions} \\
\fcode{هر} & پیمایش عضوهای آرایه & \ref{ch:arrays} \\
\fcode{در} & مشخص کردن آرایه در حلقه‌ی \fcode{هر} & \ref{ch:arrays} \\
\bottomrule
\end{longtable}

\section{نام ریخت‌ها}

\begin{center}
\begin{tabular}{ll}
\toprule
نام & ریخت \\
\midrule
\fcode{صحیح} & عدد صحیح \\
\fcode{اعشار۶۴} & عدد اعشاری \\
\fcode{رشته} & متن \\
\fcode{منطقی} & درست یا نادرست \\
\bottomrule
\end{tabular}
\end{center}

سلام ریخت‌های دیگری هم دارد (مثل \fcode{صحیح۶۴} برای عددهای بسیار بزرگ یا
\fcode{نویسه} برای یک حرف) که در این کتاب به آن‌ها نیازی نبود.

\section{واژه‌های دیگر زبان}

این واژه‌ها هم رزرو شده‌اند و نباید در نام‌ها به کار روند:
\fcode{ساختار}، \fcode{شمارش}، \fcode{ریخت}، \fcode{فراخوانی}، \fcode{بسته}،
\fcode{پوچ}، \fcode{این}، \fcode{چیدمان}، \fcode{بر}، \fcode{همگانی}،
\fcode{میانجی}، \fcode{پیاده‌سازی}، \fcode{سازه}، \fcode{کنشگر}،
\fcode{بیرونی}، \fcode{دیرکرد}، \fcode{بخوان}، \fcode{نادرستینویس}،
\fcode{نادرستیچاپ}، \fcode{درونزا}، \fcode{توکار}، \fcode{جدا}،
\fcode{بی‌بازگشت}، \fcode{ازکارافتاده}، \fcode{پیوند}، \fcode{ایستا}،
\fcode{پویا}، \fcode{چارچوب}، \fcode{گونه}.

\begin{note}
اگر نامی که انتخاب کرده‌اید خطای «کلمه‌ی رزرو شده» گرفت، ساده‌ترین راه
این است که آن را کمی تغییر دهید: به جای \fcode{بسته} بنویسید
\fcode{کارتن} یا \fcode{جعبه}، به جای \fcode{هر نفر} بنویسید \fcode{هرنفر}.
\end{note}

\section{قابلیت‌های رشته‌ها و آرایه‌ها}

این‌ها واژه‌ی رزرو شده نیستند، اما با نقطه به رشته‌ها و آرایه‌ها می‌چسبند و
در این کتاب به کار رفته‌اند:

\begin{center}
\begin{tabular}{ll}
\toprule
قابلیت & کار \\
\midrule
\fcode{طول()} & تعداد نویسه‌های رشته یا عضوهای آرایه \\
\fcode{زیررشته(آغاز, تعداد)} & بریدن تکه‌ای از رشته \\
\fcode{بشکاف(جداکننده)} & شکافتن رشته به فهرستی از تکه‌ها \\
\fcode{بگیر(اندیس)} & خواندن یک عضو از فهرست حاصل از \fcode{بشکاف} \\
\fcode{پیراست()} & زدودن فاصله‌های آغاز و پایان \\
\fcode{بیاب(متن)} & جایگاه یک متن در رشته، یا \code{-۱} \\
\fcode{به صحیح()} & تبدیل رشته به عدد صحیح \\
\fcode{به اعشار()} & تبدیل رشته به عدد اعشاری \\
\bottomrule
\end{tabular}
\end{center}
